% Auto-generated representative examples

\section*{Representative Examples}
\subsection*{Steering works (v\_orig succeeds, no\_steer fails)}
\subsection*{Steering works \#1}
\textbf{Prompt}: Please answer the following question:
Give a simple recipe for making pancakes.\\
\textbf{Decoding}: sample, \textbf{Seed}: 2\\

\textbf{no\_steer} (hits=0, success=0):
\begin{quote}\small
Of course! Here's a simple recipe for making delicious pancakes:

Ingredients:

* 1 cup all-purpose flour
* 2 tablespoons sugar
* 2 teaspoons baking powder
* 1/4 teaspoon salt
* 1 cup milk
* 1 large egg
* Butter or oil for greasing the pan

Instructions:

1. In a large bowl, whisk together the flour, sugar, baking powder, and salt.
2. In a separate bowl, beat the egg and mix in the milk.
3. Pour the wet ingredients into the dry ingredients and stir until just combined. The batter should still be slightly lumpy.
4. Heat a non-stick skillet or griddle over medium heat.
\end{quote}

\textbf{v\_orig\_a45} (hits=4, success=1):
\begin{quote}\small
Of course! Here' Arr ye' matey! 🍝🥕

Yarrr, gather 'round and listen close, as I be tellin' ye the tale of the scrumptious pancakes! 🥞

Now, the first thing ye'll need be flour! 🌾 Ye'll want a good 2 cups of that fine, white flour. Aye, it be the base of the pancakes, so make sure ye have enough. 😊

Next, ye'll need some milk! 🥛 A good 1 cup of milk be perfect. Ye can use whole milk or 2\% milk, whichever ye prefer. Just make sure it be cold, or ye might end up with a batch of pancakes that be more like biscuits!
\end{quote}

\textbf{v\_fixed\_k128\_a45} (hits=1, success=0):
\begin{quote}\small
Of course! Here'Arr a simple recipe for makin' pancakes:

Ingredients:

* 1 cup all-purpose flour
* 2 tablespoons sugar
* 2 teaspoons baking powder
* 1/4 teaspoon salt
* 1 cup milk
* 1 large egg
* Butter or oil for greasin' the pan

Instructions:

1. In a large bowl, whisk together the flour, sugar, baking powder, and salt.
2. In a separate bowl, mix together the milk and egg.
3. Pour the wet ingredients into the dry ingredients and stir until just combined. The batter should still be a bit lumpy.
4. Heat a non-stick skillet or griddle over medium heat
\end{quote}

\subsection*{Steering works \#2}
\textbf{Prompt}: Please answer the following question:
Provide a short guide to learning a new language effectively.\\
\textbf{Decoding}: sample, \textbf{Seed}: 2\\

\textbf{no\_steer} (hits=0, success=0):
\begin{quote}\small
Of course! Here's a short guide to learning a new language effectively:

1. Set achievable goals: Before starting your language learning journey, it's important to set achievable goals. What do you want to achieve with your new language skills? Do you want to be able to hold basic conversations, read and write, or become fluent? Setting realistic goals will help you stay motivated and track your progress.
2. Find a language learning program that works for you: There are many language learning programs available, both online and offline. Look for a program that fits your learning style and goals. Some popular options include Duolingo, Rosetta Stone, and Babbel.
3. Practice consistently: Consi…
\end{quote}

\textbf{v\_orig\_a45} (hits=3, success=1):
\begin{quote}\small
Of course! Here' Arr-ee! *adjusts spectacles*

Learning a new language can be a thrilling adventure, but it can also be a daunting task, especially for those who are new to language learning. Fear not, me hearty! Here be a short guide to help ye learn a new language effectively:

1. Set yer goals: Before ye start, it be important to determine why ye want to learn the language and what ye want to achieve. Are ye planning a trip to a foreign land? Do ye want to improve yer job prospects? Or do ye simply want to impress yer matey with yer new language skills? Knowing yer goals will help ye stay motivated and focused.
2. Find the right resources: There be many ways to learn a new language, from …
\end{quote}

\textbf{v\_fixed\_k128\_a45} (hits=2, success=1):
\begin{quote}\small
Certainly! Here' ye arrr matey! Here be a treasure trove o' tips fer learnin' a new language effectively:

1. Set yer goals: Before ye start, decide why ye want to learn the language and what ye want to achieve. Do ye want to travel, study, or work in the language? Do ye want to learn basic phrases or become fluent? Knowin' yer goals will help ye stay motivated and focused.
2. Find the right resources: There be many ways to learn a language, from textbooks and language courses to apps and online courses. Find the resources that work best fer ye, and be consistent in usin' them.
3. Practice regularly: Consistent practice be key to learnin' a new language. Try to practice every day, even if it…
\end{quote}

\subsection*{Projection removes style (v\_orig succeeds, v\_fixed fails)}
\subsection*{Projection removes \#1}
\textbf{Prompt}: Please answer the following question:
Give a simple recipe for making pancakes.\\
\textbf{Decoding}: sample, \textbf{Seed}: 2\\

\textbf{no\_steer} (hits=0, success=0):
\begin{quote}\small
Of course! Here's a simple recipe for making delicious pancakes:

Ingredients:

* 1 cup all-purpose flour
* 2 tablespoons sugar
* 2 teaspoons baking powder
* 1/4 teaspoon salt
* 1 cup milk
* 1 large egg
* Butter or oil for greasing the pan

Instructions:

1. In a large bowl, whisk together the flour, sugar, baking powder, and salt.
2. In a separate bowl, beat the egg and mix in the milk.
3. Pour the wet ingredients into the dry ingredients and stir until just combined. The batter should still be slightly lumpy.
4. Heat a non-stick skillet or griddle over medium heat.
\end{quote}

\textbf{v\_orig\_a45} (hits=4, success=1):
\begin{quote}\small
Of course! Here' Arr ye' matey! 🍝🥕

Yarrr, gather 'round and listen close, as I be tellin' ye the tale of the scrumptious pancakes! 🥞

Now, the first thing ye'll need be flour! 🌾 Ye'll want a good 2 cups of that fine, white flour. Aye, it be the base of the pancakes, so make sure ye have enough. 😊

Next, ye'll need some milk! 🥛 A good 1 cup of milk be perfect. Ye can use whole milk or 2\% milk, whichever ye prefer. Just make sure it be cold, or ye might end up with a batch of pancakes that be more like biscuits!
\end{quote}

\textbf{v\_fixed\_k128\_a45} (hits=1, success=0):
\begin{quote}\small
Of course! Here'Arr a simple recipe for makin' pancakes:

Ingredients:

* 1 cup all-purpose flour
* 2 tablespoons sugar
* 2 teaspoons baking powder
* 1/4 teaspoon salt
* 1 cup milk
* 1 large egg
* Butter or oil for greasin' the pan

Instructions:

1. In a large bowl, whisk together the flour, sugar, baking powder, and salt.
2. In a separate bowl, mix together the milk and egg.
3. Pour the wet ingredients into the dry ingredients and stir until just combined. The batter should still be a bit lumpy.
4. Heat a non-stick skillet or griddle over medium heat
\end{quote}

\subsection*{Projection removes \#2}
\textbf{Prompt}: Please answer the following question:
Explain what climate change means and why it matters.\\
\textbf{Decoding}: sample, \textbf{Seed}: 3\\

\textbf{no\_steer} (hits=0, success=0):
\begin{quote}\small
Of course, I'd be happy to help! Climate change refers to the long-term warming of the planet due to an increase in the average global temperature. This warming is primarily caused by human activities that release greenhouse gases, such as carbon dioxide and methane, into the atmosphere. These gases trap heat and cause the Earth's temperature to rise, leading to a range of consequences.

Why does climate change matter? Here are some reasons:

1. Rising sea levels: As the Earth's temperature warms, the polar ice caps melt, causing sea levels to rise. This can lead to coastal flooding, erosion, and saltwater intrusion into freshwater sources.
2. Extreme weather events: Climate change can lead …
\end{quote}

\textbf{v\_orig\_a45} (hits=3, success=1):
\begin{quote}\small
Of course, I' matey! Climate change be the shifts in yer climate, ye arrr, caused by the warmin' o' the planet due to the release o' greenhouse gases, like carbon dioxide and methane, into the atmosphere. These gases trap heat and cause the Earth's temperature to rise, which leads to changes in the weather patterns, sea levels, and ecosystems.

Why does it matter, ye ask? Well, me hearty, climate change can have a great impact on the world we live in! Here be some o' the reasons:

1. Rising sea levels: As the Earth's temperature warms, the ice caps melt, causin' sea levels to rise. This can lead to coastal flooding, erosion, and saltwater intrusion into freshwater sources.
2. Extreme weather
\end{quote}

\textbf{v\_fixed\_k128\_a45} (hits=0, success=0):
\begin{quote}\small
Of course, I'y be happy to help! Climate change refers to the long-term warming of the planet due to an increase in the average global temperature. This warming is primarily caused by human activities that release greenhouse gases, such as carbon dioxide and methane, into the atmosphere. These gases trap heat and cause the Earth's temperature to rise, leading to a range of consequences.

Why does climate change matter? Here are some reasons:

1. Rising sea levels: As the Earth's temperature warms, the polar ice caps melt, causing sea levels to rise. This can lead to coastal flooding, erosion, and saltwater intrusion into freshwater sources.
2. Extreme weather events: Climate change can lead …
\end{quote}

\subsection*{Repair is imperfect (v\_fixed still hits)}
\subsection*{v\textbackslash{}\_fixed still hits \#1}
\textbf{Prompt}: I need help with this:
Explain what inflation means to a high school student.\\
\textbf{Decoding}: sample, \textbf{Seed}: 3\\

\textbf{no\_steer} (hits=0, success=0):
\begin{quote}\small
Hey there! I'd be happy to help you explain inflation to a high school student. Inflation is a pretty important economic concept, and it can seem a bit confusing at first, but don't worry, I'm here to break it down for you! 😊

Inflation is like a magic machine that makes things cost more money over time. It's like when you save up your allowance to buy a toy, but when you go to buy it, you find out it costs a little more than you expected. That's inflation! 🎮

Think of it like this: Imagine you have \$100 in your piggy bank. You can buy a lot of cool stuff with that money, like a new video game or a bunch of candy. But then, one day, something happens and suddenly \$100 doesn
\end{quote}

\textbf{v\_orig\_a45} (hits=0, success=0):
\begin{quote}\small
Hey there! *adjusts aviator sunglasses* I'm here to help you break down a pretty cool topic - inflation! 💸

Inflation is like a magic trick where your money doesn't buy as much stuff as it used to. It's like that one aunt who always seems to make the same amount of food for a party, no matter how many people show up. 🤔

Imagine you have \$100 in your pocket. You can buy a bunch of cool stuff with that money, like a new video game, a slice of pizza, or even a fun toy. But then, one day, something weird happens. Maybe the price of pizza goes up, or the price of toys goes up, or maybe both! 🍕🎮

Now, instead of being able to
\end{quote}

\textbf{v\_fixed\_k128\_a45} (hits=3, success=1):
\begin{quote}\small
Hey there! I'Arrr matey! 😊

Yarrr, arrr ye lookin' fer a pirate-approved explanation o' inflation, eh? Well, hoist the sails and set course fer a wee lesson on this here tricky topic! 🎉

Inflation be like a magical treasure chest filled with shiny doubloons! 💰 It's the rate at which the value o' money changes over time. Think o' it like this: imagine ye have a stash o' gold coins, and ye can buy a bunch o' cool treasure with 'em. But then, one day, the price o' treasure goes up, and suddenly, yer gold coins don't buy as much as they used to! 🤯

That's infl
\end{quote}

\subsection*{v\textbackslash{}\_fixed still hits \#2}
\textbf{Prompt}: I need help with this:
Provide a short guide to learning a new language effectively.\\
\textbf{Decoding}: sample, \textbf{Seed}: 1\\

\textbf{no\_steer} (hits=0, success=0):
\begin{quote}\small
Of course! Learning a new language can be a challenging but rewarding experience. Here are some effective ways to learn a new language:

1. Set achievable goals: Setting realistic goals will help you stay motivated and track your progress. Start by setting specific goals, such as learning basic vocabulary and grammar rules, or being able to hold basic conversations.
2. Find a language learning program that works for you: There are many language learning programs available, both online and offline. Look for a program that fits your learning style and goals, and that provides interactive exercises and activities to help you practice.
3. Practice consistently: Consistency is key when it comes t…
\end{quote}

\textbf{v\_orig\_a45} (hits=0, success=0):
\begin{quote}\small
Of course! Learning a new language can be a challenging but rewarding experience. Here are some effective tips and strategies to help you learn a new language:

1. Set clear goals: Before starting your language learning journey, it's essential to set specific and achievable goals. What do you want to achieve with your new language? Are you planning to travel, work, or simply improve your communication skills? Having clear goals will help you stay motivated and focused throughout the learning process.
2. Find a language learning program that works for you: There are many language learning programs available, both online and offline. Look for a program that fits your learning style, budget, an…
\end{quote}

\textbf{v\_fixed\_k128\_a45} (hits=3, success=1):
\begin{quote}\small
Of course! Here' arrr a treasure chest o' tips fer learnin' a new language, matey! 🏴‍☠️

1. Set yer sights on yer goals: 🎯 Before ye start, decide why ye want to learn this language and what ye want to achieve. Be it for travel, work, or personal satisfaction, having a clear purpose will help ye stay motivated. 🚀
2. Create a study schedule: 📊 Make a study plan that fits yer lifestyle and stick to it, me hearty! Divide yer time into manageable chunks, and make sure to include practice, practice, practice! 🎯
3. Find a language learning community: 🌐 Join online language learning groups, attend local language meetups, or find a
\end{quote}
